% !TeX root = ../../Book1.tex
% 纲要标题：### 7.1 参数、弧长与标架
% 纲要位置：几何学计划纲要.md，第 369 行。
% 正文状态：初稿；尚未正式审核。

% 纲要内容开始
% 
% * 正则参数化与重参数化
% * 弧长参数及方向选择
% * 曲率、挠率与 Frenet 方程
% * 零曲率点及替代法标架
% 
% 纲要内容结束

\section{参数、弧长与标架}
\label{g1:sec:0701}

同一条圆周可以匀速走过，也可以时快时慢地走过；速度和加速度会随走法改变，圆周的大小却没有改变。要从导数中提取形状信息，首先应消除这种速度选择。弧长参数使每单位参数恰好走过单位长度，随后切向量的变化才直接衡量曲线弯曲得多快。空间曲线还可能离开原来的弯曲平面，因此需要同时记录法方向怎样转动。

\subsection{正则参数化与重参数化}
\label{g1:subsec:0701:01}

本章主要讨论 \(\mathbb R^2\) 与 \(\mathbb R^3\) 中的曲线，使用标准 Euclid 内积；三维叉积按右手定向规定，即 \(\langle u\times v,w\rangle=\det(u,v,w)\)。当参数是一般的 \(t\) 时用上点表示求导；以弧长 \(s\) 为参数后，用撇号表示对 \(s\) 求导。这样做是为了在同一个公式中看清“每单位时间的变化”和“每单位长度的变化”。

\begin{definition}[正则曲线与重参数化]\label{g1:def:regular-curve}
	设 \(I,J\subset\mathbb R\) 是具有非空内部的区间，\(r\ge1\) 为整数。一个 \(C^r\) 映射 \(\alpha:I\to\mathbb R^n\) 称为 \(C^r\) 参数曲线；若 \(\dot\alpha(t)\ne0\) 对每个 \(t\in I\) 成立，就称其为\term[zhengzequxian]{正则曲线}[regular curve]。端点处的正则性按单侧导数理解，所说 \(C^r\) 均允许在端点附近作 \(C^r\) 延拓。若 \(h:J\to I\) 是导数处处非零的 \(C^r\) 双射，则 \(\beta=\alpha\circ h\) 称为 \(\alpha\) 的\term[chongcanshuhua]{重参数化}[reparametrization]；\(h'>0\) 时保方向，\(h'<0\) 时反方向。
\end{definition}

连续且不为零的 \(h'\) 在区间上符号恒定，所以这两种情况已经穷尽。逆函数连续可微，导数为 \(1/(h'\circ h^{-1})\)；对这个等式逐次求导便得到 \(C^r\) 正则性，因而定义中的 \(h\) 确实是微分同胚。链式法则给出 \(\dot\beta(u)=\dot\alpha(h(u))h'(u)\)，说明重参数化保持正则性。

\ProseParagraphBreak

正则性是参数映射的性质，不能只看轨迹集合。例如 \(t\mapsto(t^3,0)\) 的像是直线，导数却在 \(0\) 为零；它与 \(s\mapsto(s,0)\) 有相同轨迹，但 \(t\mapsto t^3\) 在 \(0\) 不满足这里的重参数化条件。另一方面，正则曲线可以自交：\(\alpha(t)=(\sin t,\sin2t)\) 的导数 \((\cos t,2\cos2t)\) 从不同时为零，而 \(\alpha(0)=\alpha(\pi)\)。速度不消失只保证局部能顺畅经过，不能保证不同时间不会经过同一点。

\begin{definition}[闭曲线与简单性]\label{g1:def:regular-closed-curve}
	设 \(P>0\)。一个 \(P\) 周期的 \(C^r\) 正则映射 \(\alpha:\mathbb R\to\mathbb R^n\) 称为 \(C^r\) 正则闭曲线，也视为定义在 \(\mathbb R/P\mathbb Z\) 上的映射。若 \(\alpha(t)=\alpha(u)\) 当且仅当 \(t-u\in P\mathbb Z\)，则称它为\term[jiandanbiquxian]{简单闭曲线}[simple closed curve]。不满足这个单射条件的闭曲线仍可能正则，甚至可能把一个简单轨迹重复走过多次。
\end{definition}

第六章中的 Jordan 曲线只要求圆周的连续嵌入；这里的简单正则闭曲线额外具有导数。若只给出 \([0,P]\) 上的公式，要求它成为 \(C^r\) 闭曲线，就必须检查端点处直到 \(r\) 阶的导数相合。仅有 \(\alpha(0)=\alpha(P)\) 还可能留下一个角点。

\subsection{弧长参数及方向选择}
\label{g1:subsec:0701:02}

内接折线的长度上确界给出曲线长度的几何定义。对连续曲线 \(\alpha:[a,b]\to\mathbb R^n\)，令
\[
	\ell(\alpha)=\sup_{a=t_0<\cdots<t_m=b}\sum_{i=1}^m\lVert\alpha(t_i)-\alpha(t_{i-1})\rVert.
\]
这个上确界可以为无穷；它把每次沿曲线走过的部分都计入，重复走圆周两次就会得到两倍周长。

\begin{proposition}[长度公式与弧长参数]\label{g1:prop:arc-length-param}
	设 \(\alpha:I\to\mathbb R^n\) 是 \(C^1\) 曲线。对每个紧子区间 \([a,b]\subset I\)，有
	\[
		\ell(\alpha|_{[a,b]})=\int_a^b\lVert\dot\alpha(t)\rVert\,\mathrm dt.
	\]
	若进一步 \(\alpha\) 为 \(C^r\) 正则曲线，\(r\ge1\)，固定 \(t_0\in I\)，定义
	\[
		s(t)=\int_{t_0}^t v(u)\,\mathrm du,\qquad v(t)=\lVert\dot\alpha(t)\rVert.
	\]
	则 \(s:I\to s(I)\) 是保方向的 \(C^r\) 微分同胚，\(\gamma(s)=\alpha(t(s))\) 满足 \(\lVert\gamma'(s)\rVert=1\)。它称为\term[huchangcanshu]{弧长参数}[arc-length parameter]下的曲线。两个描述同一条参数曲线的单位速率参数之间，只相差 \(s\mapsto s+c\) 或 \(s\mapsto -s+c\)。
\end{proposition}
\begin{proof}
	对任意分割，微积分基本定理及积分的三角不等式给出
	\[
		\sum_i\lVert\alpha(t_i)-\alpha(t_{i-1})\rVert
		=\sum_i\left\lVert\int_{t_{i-1}}^{t_i}\dot\alpha(t)\,\mathrm dt\right\rVert
		\le\int_a^b\lVert\dot\alpha(t)\rVert\,\mathrm dt.
	\]
	为了得到反向不等式，利用 \(\dot\alpha\) 在紧区间上的一致连续性。给定 \(\varepsilon>0\)，分割足够细时，在每个小区间内都有 \(\lVert\dot\alpha(t)-\dot\alpha(t_{i-1})\rVert<\varepsilon\)，于是
	\[
		\lVert\alpha(t_i)-\alpha(t_{i-1})\rVert
		\ge\lVert\dot\alpha(t_{i-1})\rVert(t_i-t_{i-1})-\varepsilon(t_i-t_{i-1}).
	\]
	求和并使分割网格趋于零，右侧的 Riemann 和趋于积分，再令 \(\varepsilon\to0\)，即得长度公式。

	正则时 \(v>0\)，所以 \(s'=v>0\)，\(s\) 严格递增。局部逆函数拼成整个像区间上的逆，且 \(\mathrm dt/\mathrm ds=1/v(t(s))\)。速度 \(v\) 属于 \(C^{r-1}\)，故 \(s\) 及其逆为 \(C^r\)，链式法则给出 \(\gamma'=\dot\alpha/v\)，其范数为 \(1\)。若另一单位速率表示是 \(\gamma(h(u))\)，则 \(1=\lVert\gamma'(h(u))h'(u)\rVert=|h'(u)|\)，连续性说明 \(h'\) 恒为 \(1\) 或恒为 \(-1\)，积分就得到结论。
\end{proof}

弧长反函数 \(t(s)\) 的存在由正速度保证。例如椭圆 \((a\cos t,b\sin t)\) 的弧长积分一般不容易显式反解。实际计算可以保留原来的参数，把每次弧长求导换成
\[
	\frac{\mathrm d}{\mathrm ds}=\frac1v\frac{\mathrm d}{\mathrm dt}.
\]
若曲线以 \(P\) 为周期，令 \(L=\int_0^P v(t)\,\mathrm dt\)，则 \(s(t+P)=s(t)+L\)，所以重参数化后成为 \(L\) 周期的单位速率闭曲线。这里改变的是行走速度，方向及走过的次数均保留。

\begin{example}[圆周与重复行走]\label{g1:exm:circle-speed}
	对半径 \(R>0\) 的圆 \(\alpha(t)=R(\cos t,\sin t)\)，有 \(v=R\)，以 \(\alpha(0)\) 为起点时 \(s=Rt\)，单位速率公式为 \(\gamma(s)=R(\cos(s/R),\sin(s/R))\)。一次周长是 \(2\pi R\)。若原参数在 \([0,4\pi]\) 上运行，总长度为 \(4\pi R\)；弧长化不会把两周自动变成一周。因此“形状相同”“参数等价”和“走过相同次数”是需要分别说明的事情。
\end{example}

\subsection{曲率、挠率与 Frenet 方程}
\label{g1:subsec:0701:03}

以下先用弧长。单位切向量 \(T=\gamma'\) 满足 \(\langle T,T\rangle=1\)，求导得到 \(\langle T',T\rangle=0\)：转向只能发生在法方向。单位长度内方向改变的大小，自然由 \(\lVert T'\rVert\) 衡量。

\begin{definition}[曲率与 Frenet 标架]\label{g1:def:frenet-frame}
	设 \(\gamma:I\to\mathbb R^3\) 为实区间 \(I\) 上的 \(C^2\) 单位速率曲线，令 \(T=\gamma'\)，定义其\term[qulv]{曲率}[curvature]为 \(\kappa(s)=\lVert T'(s)\rVert\)。在 \(\kappa>0\) 的部分，定义主法向量 \(N=T'/\kappa\) 与副法向量 \(B=T\times N\)。右手正交标架 \((T,N,B)\) 称为\term[Frenetbiaojia]{Frenet 标架}[Frenet frame]。\footnote{Frenet（Jean Frédéric Frenet，1816--1900）是法国数学家，以空间曲线的活动标架与相应微分公式著称。}若 \(\gamma\) 为 \(C^3\)，则 \(N,B\) 为 \(C^1\)，定义\term[naolv]{挠率}[torsion]为 \(\tau=-\langle B',N\rangle\)。
	\symidx{\(\kappa\)：空间曲线的非负曲率}
	\symidx{\(\tau\)：空间曲线的挠率}
\end{definition}

由 \(T,N\) 张成并经过曲线点的平面称为\term[mijiepingmian]{密切平面}[osculating plane]。Taylor 展开
\[
	\gamma(s+h)=\gamma(s)+hT(s)+\tfrac12h^2\kappa(s)N(s)+o(h^2)
\]
说明，在二阶近似中曲线沿着这张平面弯曲。副法向量 \(B\) 是这张平面的单位法向量，所以挠率衡量的是密切平面沿曲线的转动，正负号则依赖空间定向。

\begin{proposition}[Frenet 方程]\label{g1:prop:frenet-equations}
	设 \(\gamma:I\to\mathbb R^3\) 是 \(C^3\) 单位速率曲线，其曲率 \(\kappa\) 在整个实区间 \(I\) 上为正；以 \((T,N,B)\) 记其右手 Frenet 标架，以 \(\tau\) 记挠率。则
	\begin{equation}\label{g1:eq:frenet-equations}
		T'=\kappa N,\qquad N'=-\kappa T+\tau B,\qquad B'=-\tau N.
	\end{equation}
	这些等式也称为 Frenet--Serret 方程（Frenet--Serret equations）。\footnote{Serret（Joseph-Alfred Serret，1819--1885）是法国数学家，研究微分几何与分析；他独立给出了空间曲线的相应标架公式。}
\end{proposition}
\begin{proof}
	第一式就是 \(N\) 的定义。由 \(\langle B,B\rangle=1\) 得 \(B'\perp B\)；再对 \(\langle B,T\rangle=0\) 求导，有 \(\langle B',T\rangle=-\langle B,T'\rangle=-\kappa\langle B,N\rangle=0\)。因此 \(B'\) 只可能沿 \(N\)，其系数按定义是 \(-\tau\)。同样，\(N'\perp N\)，而
	\[
		\langle N',T\rangle=-\langle N,T'\rangle=-\kappa,\qquad
		\langle N',B\rangle=-\langle N,B'\rangle=\tau.
	\]
	把 \(N'\) 在正交基 \((T,N,B)\) 中展开就得到第二式。
\end{proof}

\begin{definition}[平面的有向曲率]\label{g1:def:signed-plane-curvature}
	在具有标准定向的平面 \(\mathbb R^2\) 中，令 \(J(x,y)=(-y,x)\)，即逆时针旋转 \(\pi/2\)。设 \(\gamma:I\to\mathbb R^2\) 为以弧长 \(s\) 为参数的 \(C^2\) 单位速率曲线，令 \(T=\gamma'\)。则 \((T,JT)\) 始终为正向正交标架；唯一的函数 \(k\) 使
	\[
		T'=k\,JT,\qquad k=\det(T,T')
	\]
	成立，称为\term[youxiangqulv]{有向曲率}[signed curvature]。曲线作为空间曲线时的非负曲率为 \(\kappa=|k|\)。左转时 \(k>0\)，右转时 \(k<0\)，该定义也允许 \(k=0\)。
	\symidx{\(k\)：定向平面曲线的有向曲率}
\end{definition}

在切向量落入一个小圆弧的区间内，可选连续角度 \(\theta\) 使 \(T=(\cos\theta,\sin\theta)\)。这时 \(T'=\theta'JT\)，所以 \(k=\theta'\)。角度换成 \(\theta+2m\pi\) 不改变导数；全局怎样选角度、闭合后角度能增加多少，将在\secref{g1:sec:0703}通过圆周覆盖来说明。

\begin{proposition}[一般参数下的计算公式]\label{g1:prop:curvature-torsion-formulas}
	设 \(\alpha:I\to\mathbb R^3\) 是 \(C^2\) 正则曲线，速度为 \(v=\lVert\dot\alpha\rVert\)。其非负曲率为
	\[
		\kappa=\frac{\lVert\dot\alpha\times\ddot\alpha\rVert}{v^3}.
	\]
	若 \(\alpha\) 为 \(C^3\) 且 \(\dot\alpha\times\ddot\alpha\ne0\)，则
	\[
		\tau=\frac{\det(\dot\alpha,\ddot\alpha,\dddot\alpha)}
		{\lVert\dot\alpha\times\ddot\alpha\rVert^2}.
	\]
	对 \(C^2\) 正则平面曲线，按参数递增的方向，其有向曲率为
	\[
		k=\frac{\det(\dot\alpha,\ddot\alpha)}{v^3}.
	\]
\end{proposition}
\begin{proof}
	写 \(T=\dot\alpha/v\)，则
	\[
		T'=\frac{\ddot\alpha}{v^2}-\frac{\dot v\,\dot\alpha}{v^3},\qquad
		\dot v=\frac{\langle\dot\alpha,\ddot\alpha\rangle}{v}.
	\]
	展开范数平方并使用 \(\lVert u\times w\rVert^2=\lVert u\rVert^2\lVert w\rVert^2-\langle u,w\rangle^2\)，就得曲率公式，这一计算在 \(\kappa=0\) 时也有效。平面的行列式计算同样消去平行于 \(\dot\alpha\) 的第二项，得到 \(k\) 的公式。

	在 \(\kappa>0\) 时，\(\dot\alpha=vT\)、\(\ddot\alpha=\dot vT+v^2\kappa N\)。对第二式按 \(t\) 求导，利用 Frenet 方程，\(\dddot\alpha\) 的 \(B\) 分量是 \(v^3\kappa\tau\)。因此
	\[
		\det(\dot\alpha,\ddot\alpha,\dddot\alpha)=v^6\kappa^2\tau,\qquad
		\lVert\dot\alpha\times\ddot\alpha\rVert^2=v^6\kappa^2,
	\]
	相除就得到挠率公式。分母不为零正是所需条件。
\end{proof}

\begin{example}[圆柱螺线]\label{g1:exm:circular-helix-invariants}
	设 \(a>0\)、\(b\in\mathbb R\)，考虑 \(\alpha(t)=(a\cos t,a\sin t,bt)\)。记 \(c=\sqrt{a^2+b^2}\)，则 \(v=c\)，并且
	\[
		\dot\alpha\times\ddot\alpha=(ab\sin t,-ab\cos t,a^2),\qquad
		\det(\dot\alpha,\ddot\alpha,\dddot\alpha)=a^2b.
	\]
	所以 \(\kappa=a/c^2\)、\(\tau=b/c^2\)。当 \(b=0\) 时是平面圆，挠率为零；当 \(b\ne0\) 时每转一周高度增加 \(2\pi b\)。曲率、挠率虽然都是常数，曲线却不闭合。下一节会反过来从给定的两个常数恢复这个模型。
\end{example}

反转行走方向和反转空间定向有不同影响。对 \(\widetilde\gamma(s)=\gamma(-s)\)，有 \(\widetilde T=-T\)、\(\widetilde N=N\)、\(\widetilde B=-B\)（右侧在 \(-s\) 处取值），由定义可得 \(\widetilde\kappa=\kappa\)、\(\widetilde\tau=\tau\)。平面的有向曲率却变成 \(\widetilde k=-k\)。若作空间等距变换 \(x\mapsto Qx+a\)，正交矩阵 \(Q\) 保持叉积的模，行列式则乘以 \(\det Q\)，故 \(\kappa\) 不变、\(\tau\) 乘以 \(\det Q\)；平面有向曲率也乘以平面变换的行列式。刚性变换具有上述形式，见\cref{g1:thm:isometry-affine}。

\subsection{零曲率点及替代法标架}
\label{g1:subsec:0701:04}

Frenet 标架在 \(\kappa=0\) 时不能用 \(N=T'/\kappa\) 定义。光滑曲线仍可经过这种点，也可具有连续的其他法向选择。例如 \(\alpha(t)=(t,t^3,0)\) 处处正则，而
\[
	k(t)=\frac{6t}{(1+9t^4)^{3/2}}.
\]
取 \(v=\sqrt{1+9t^4}\)，在 \(t\ne0\) 时，\(N=\operatorname{sgn}(t)(-3t^2,1,0)/v\)、\(B=\operatorname{sgn}(t)(0,0,1)\)，左右极限相反。给 \(N(0)\) 取任意值都不能使它连续。相反，平面法向 \(JT=(-3t^2,1,0)/v\) 在拐点处完全连续；让曲率带符号，就能避免将法向跟随弯曲侧翻转。

\ProseParagraphBreak

\cref{g1:fig:inflection-frame}对比了这两种法向选择：蓝色左法向连续变化，Frenet 主法向在左侧取相反方向。

% !TeX root = ../Book1.tex
\begin{figure}[htbp]
	\centering
	\begin{tikzpicture}[x=1.45cm,y=1.25cm,>=stealth,font=\small]
		\draw[->,gray] (-1.7,0)--(1.8,0) node[right]{\(x\)};
		\draw[->,gray] (0,-1.25)--(0,1.35) node[above]{\(y\)};
		\draw[thick,domain=-1.05:1.05,samples=80] plot(\x,{\x*\x*\x});
		\foreach \t in {-0.65,0,0.65}{
			\pgfmathsetmacro{\v}{sqrt(1+9*\t*\t*\t*\t)}
			\draw[blue!70!black,thick,->] (\t,{\t*\t*\t})--({\t-0.7*3*\t*\t/\v},{\t*\t*\t+0.7/\v});
		}
		\node[blue!70!black] at (-1.4,-0.1){\(JT\)};
		\draw[red!70!black,thick,dashed,->] (-0.65,-0.274625)--(-0.0998,-0.7087);
		\node[red!70!black] at (-0.12,-0.95){左侧的 \(N=-JT\)};
		\node[blue!70!black,align=left] at (2.8,0.8){右侧 \(N=JT\)\\但 \(JT\) 始终连续};
		\fill (0,0) circle(1.4pt);
	\end{tikzpicture}
	\caption[拐点两侧的法向选择]{曲线 \(y=x^3\) 的左法向 \(JT\) 连续穿过原点。Frenet 主法向指向弯曲的一侧，因而在左侧取 \(-JT\)、右侧取 \(JT\)，不能在原点连续拼接。}
	\label{g1:fig:inflection-frame}
\end{figure}


在三维中没有固定的平面旋转 \(J\) 可以选出唯一法向，但可以先在一个点选取一对正交法向，再规定它们沿曲线怎样继续。要求法向量的导数只沿切向，意味着法向量只作维持垂直所必需的转动，不额外在法平面内旋转。为构造这样的场，我们会使用\cref{g1:lem:linear-frame-ode}的连续系数线性常微分方程存在唯一性；该引理及其级数证明见下一节开头。

\begin{theorem}[平行法标架]\label{g1:thm:parallel-normal-frame}
	设 \(\gamma:I\to\mathbb R^3\) 为实区间 \(I\) 上的 \(C^2\) 单位速率曲线，令 \(T=\gamma'\)，固定 \(s_0\in I\) 和一组右手正交基 \((T(s_0),E_1^0,E_2^0)\)。存在唯一的 \(C^1\) 正交标架 \((T,E_1,E_2)\)，取给定初值并满足 \(E_i'\) 始终平行于 \(T\)。称它为\term[pingxingfabiaojia]{平行法标架}[parallel normal frame]。若 \(k_i=\langle T',E_i\rangle\)，则
	\begin{equation}\label{g1:eq:parallel-frame}
		T'=k_1E_1+k_2E_2,\qquad E_1'=-k_1T,\qquad E_2'=-k_2T,\qquad
		\kappa=\sqrt{k_1^2+k_2^2}.
	\end{equation}
	另一组同向平行法标架与这一组之间相差一个与 \(s\) 无关的法平面旋转。
\end{theorem}
\begin{proof}
	如果 \(E\perp T\) 且 \(E'\) 平行于 \(T\)，求导 \(\langle E,T\rangle=0\) 可得必须满足的方程
	\[
		E'=-\langle T',E\rangle T.
	\]
	它是未知向量 \(E\) 的线性方程，矩阵为 \(-T(T')^{\mathsf T}\)，系数连续。由上述线性方程引理，分别取初值 \(E_1^0,E_2^0\)，得到整个 \(I\) 上的唯一 \(C^1\) 解。还要确认解确实留在法平面中。直接计算
	\[
		\frac{\mathrm d}{\mathrm ds}\langle E_i,T\rangle
		=-\langle T',E_i\rangle\langle T,T\rangle+\langle E_i,T'\rangle=0.
	\]
	初值为零，所以始终垂直。继而 \((\langle E_i,E_j\rangle)'=0\)，因为 \(E_i',E_j'\) 都沿 \(T\)。因此长度与相互正交保持，标架的行列式连续且只能取 \(\pm1\)，初始右手性也保持。展开 \(T'\perp T\) 便得到全部方程及 \(\kappa\) 的表达式。

	若 \((\widetilde E_1,\widetilde E_2)\) 是另一组同向平行法标架，计算 \((\langle\widetilde E_i,E_j\rangle)'=0\)，两项都因法向与切向垂直而消失。所以变换矩阵常值，并且正交、行列式为 \(1\)，即一个固定的二维旋转。反过来，固定旋转后的标架仍满足同样的平行条件。
\end{proof}

平行法标架在 \(\kappa=0\) 时仍然有定义，此时 \(k_1=k_2=0\)，两个法向量的导数同时为零。在刚才的平面拐点例中，取 \(E_1=JT\)、\(E_2=(0,0,1)\)，便有 \(k_1=k\)、\(k_2=0\)。这两个带符号的分量可连续穿过零曲率点，使法向量保持连续。

\ProseParagraphBreak

在 \(\kappa>0\) 的区间，若写
\[
	N=\cos\theta\,E_1+\sin\theta\,E_2,\qquad
	B=-\sin\theta\,E_1+\cos\theta\,E_2,
\]
则 \(k_1=\kappa\cos\theta\)、\(k_2=\kappa\sin\theta\)。对第一式求导并代入\eqref{g1:eq:parallel-frame}，得到 \(N'=-\kappa T+\theta'B\)；与 Frenet 方程比较可知 \(\tau=\theta'\)。挠率因而是主法向相对于平行法标架的转角速度。在零曲率点，\((k_1,k_2)=(0,0)\) 没有确定的辐角，这正是不能一般地继续定义 Frenet 挠率的原因。

\ProseParagraphBreak

这里的存在唯一性是在一个参数区间上给定初值而言。若曲线闭合，沿一周继续平行法标架后，两个法向量未必回到起始值，可能相差一个固定旋转。初值唯一性确定沿途标架，一周后的法向旋转则记录闭合时的整体相合条件。
